\section{2024-11-13 Klausur}

\startplacetable[location={here},title={Noten}]
\startuseMPgraphic{table:incell:rarrow}
	drawarrow (1mm,0)--(\overlaywidth - 1mm,0);
\stopuseMPgraphic
\defineoverlay
  [tableIncellRarrow]
  [\useMPgraphic{table:incell:rarrow}]
\setupTABLE[each][offset=1mm,frame=off]
\setupTABLE[c][1][background=color,backgroundcolor=gray,rightframe=on]
\setupTABLE[r][2][topframe=on]
\bTABLE
\startCSV
ab, 20\%, 27\%, 33\%, 40\%, 45\%, 5\%-Schritte ..., 60\%, 80\%, 90\%
Note, 01, 02, 03, 04, 05, [background=tableIncellRarrow], 08, 12, 14
\stopCSV
\eTABLE
\stopplacetable

ggf. Abzüge Sprachrichtigkeit

\startplacetable[location={here},title={Abzüge Sprachrichtigkeit}]
\setupTABLE[each][offset=1mm,frame=off]
\setupTABLE[c][1][background=color,backgroundcolor=gray,rightframe=on]
\setupTABLE[r][2][topframe=on]
\bTABLE
\startCSV
5F/1S, -01P
7F/1S, -02P
\stopCSV
\eTABLE
\stopplacetable

\startframedtipp
Für Anki zum Lernen

\startitemize[packed]
\item Trigonomie
\startitemize[packed]
\item sinus
\item arctan
\item etc
\stopitemize
\item 10er Potenzen, Präfixe, Vorsilben
\stopitemize
\stopframedtipp

\section{2024-11-13 Doppelspalt Messung}

Wellenlänge-Spannungs-Diag
